\chapter{Alcance y Decisiones Metodológicas}

\section{Alcance analítico}

La implementación observada en el repositorio trabaja con una unidad de análisis comunal: una fila del dataset final corresponde a una comuna. El universo final se restringe a la Provincia de Santiago y se valida contra \texttt{data/raw/dim\_comuna\_base.csv}, que contiene 32 comunas. El producto integrado resultante mantiene esa cobertura sin faltantes ni extras.

El análisis tiene naturaleza comparativa y descriptiva. En consecuencia, los outputs de Fase 9 se interpretan como patrones observados entre comunas y no como evidencia para inferir relaciones causales. Este criterio aparece de manera consistente en la documentación del repositorio y en los reportes analíticos.

\section{Corte temporal y consistencia de referencia}

El proyecto no fuerza una homogeneización artificial de años entre fuentes. En el dataset final quedan registrados explícitamente \texttt{anio\_poblacion = 2024}, \texttt{anio\_pobreza = 2022}, \texttt{anio\_areas\_verdes = 2024} y \texttt{anio\_ingresos = 2024}. Esta decisión preserva trazabilidad respecto del origen de cada variable y evita ocultar diferencias temporales entre insumos.

\section{Llave principal y producto final}

La llave principal obligatoria del proyecto es \texttt{codigo\_comuna}. \texttt{nombre\_comuna} se mantiene como apoyo descriptivo y de verificación, pero no se usa como llave de integración. El producto final observado se compone de un CSV integrado con 32 filas y 12 columnas, una base SQLite con tres tablas y un conjunto de outputs de validación y análisis exploratorio.

\section{Síntesis de decisiones}

\begin{table}[htbp]
\centering
\footnotesize
\caption{Decisiones metodológicas observadas en el repositorio.}
\begin{tabular}{|p{3.0cm}|p{5.6cm}|p{5.0cm}|}
\hline
Dimensión & Decisión observada & Evidencia principal \\
\hline
Unidad de análisis & Una fila = una comuna & README.md; outputs/resumen\_dataset\_final.md \\
Cobertura territorial & Provincia de Santiago, 32 comunas & data/raw/dim\_comuna\_base.csv; outputs/reporte\_validacion\_final.md \\
Llave principal & codigo\_comuna & outputs/resumen\_homologacion.md; src/comunas.py \\
Naturaleza analítica & Comparativa y descriptiva, no causal & README.md; outputs/analisis\_exploratorio.md \\
Temporalidad & Pobreza 2022; población, áreas verdes e IPP 2024 & data/raw/metadata\_fuentes.csv; outputs/resumen\_dataset\_final.md \\
Integración & Merges left con validación one\_to\_one sobre codigo\_comuna & outputs/log\_integracion.md; src/integrate.py \\
Tratamiento de faltantes & Sin imputación; métricas derivadas quedan NA solo si la división es inválida & outputs/resumen\_dataset\_final.md; src/integrate.py \\
Producto final & CSV final, SQLite y outputs analíticos reproducibles & README.md; outputs/reporte\_carga\_sqlite.md \\
\hline
\end{tabular}
\end{table}

